% 图 4.3 TVM 优化逻辑流程图 — 基于 CICC0903540 技术文档 4.2.1 与 4.2.3
\documentclass[tikz,border=8pt]{standalone}
% 顶会/顶刊风格：低饱和度马卡龙 + 高级感
% 适用于 NeurIPS, CVPR, ICLR, IEEE TPAMI
\usepackage{amssymb}
\usepackage{fontspec}
\usepackage{xeCJK}
\setCJKmainfont{SimHei}
\usepackage{tikz}
\usepackage{pgfplots}
\usetikzlibrary{
  arrows.meta, positioning, shapes.geometric, calc,
  backgrounds, fit, decorations.pathreplacing, patterns
}

% 低饱和度马卡龙配色
\definecolor{mBlue}{HTML}{AEC6CF}
\definecolor{mPink}{HTML}{FFD1DC}
\definecolor{mGreen}{HTML}{B1D8B7}
\definecolor{mPurple}{HTML}{B39EB5}
\definecolor{mYellow}{HTML}{FDFD96}
\definecolor{mGray}{HTML}{E5E4E2}
\definecolor{mDark}{HTML}{4A4A4A}

% 全局 TikZ 样式
\tikzset{
  >=stealth,
  every node/.style={align=center, font=\sffamily},
  block/.style={
    rounded corners=3pt, draw=mDark, align=center,
    minimum height=0.8cm, inner sep=5pt},
  arr/.style={-{Stealth[length=4pt]}, semithick, draw=mDark},
  % 信息密度增强：张量维度标注
  ts/.style={font=\tiny\sffamily, color=mDark!80},
  % 数学操作符节点（⊕加 ⊗乘 ‖拼接）
  op/.style={circle, draw=mDark, minimum size=0.4cm, inner sep=0,
    font=\scriptsize\sffamily, fill=mGray!60},
  % 图例
  legendbox/.style={rectangle, draw=mDark!50, rounded corners=2pt,
    fill=mGray!15, inner sep=4pt, font=\tiny\sffamily},
}

\begin{document}
\begin{tikzpicture}[
  node distance=0.6cm and 0.6cm,
  stage/.style={
    rectangle, draw=mDark, fill=mBlue!40,
    minimum width=1.9cm, minimum height=0.6cm,
    font=\scriptsize\sffamily, align=center, rounded corners=3pt,
    semithick, inner sep=6pt},
  arrow/.style={-{Stealth[length=4pt]}, semithick, draw=mDark},
  lbl/.style={font=\scriptsize\sffamily\bfseries, color=mDark},
]
  \node[lbl] (t1) {1};
  \node[stage, right=0.25cm of t1] (b1) {compat};
  \node[stage, right=0.6cm of b1] (c1) {safe runtime};
  \draw[arrow] (b1) -- (c1);

  \node[lbl, below=0.6cm of t1] (t2) {2};
  \node[stage, right=0.25cm of t2] (b2) {generic};
  \node[stage, right=0.6cm of b2] (c2) {cortex-a72};
  \draw[arrow] (b2) -- (c2);

  \node[lbl, below=0.6cm of t2] (t3) {3};
  \node[stage, right=0.25cm of t3] (b3) {rebuild};
  \node[stage, right=0.6cm of b3] (c3) {增量调优};
  \draw[arrow] (b3) -- (c3);

  \node[lbl, below=0.6cm of t3] (t4) {4};
  \node[stage, right=0.25cm of t4] (b4) {payload};
  \node[stage, right=0.6cm of b4] (c4) {真实重建};
  \draw[arrow] (b4) -- (c4);

  \node[lbl, below=0.6cm of t4] (t5) {5};
  \node[stage, right=0.25cm of t5] (b5) {人工};
  \node[stage, right=0.6cm of b5] (c5) {SHA guard};
  \draw[arrow] (b5) -- (c5);

  \draw[decorate, decoration={brace, amplitude=2pt, raise=1pt}, semithick, draw=mGray]
    (b1.north west) -- (c5.south east)
    node[midway, above=8pt, font=\scriptsize\sffamily\bfseries, color=mDark]
    {baseline $\rightarrow$ current};

  % 图例（预留边距）
  \node[legendbox, anchor=north east] at ([xshift=-5pt, yshift=-5pt]current bounding box.south east)
    [align=left, text width=3cm] {
    \textbf{图例}\\[1pt]
    1 运行时 \quad 2 target \quad 3 调优\\
    4 评测 \quad 5 治理
  };
\end{tikzpicture}
\end{document}
